% Après lecture de votre rapport PFA intitulé « FakeFinder – Système Intelligent de Détection du Hameçonnage par E-mail », voici mon évaluation.
%
% Note proposée : 15 / 20
%
% Votre projet traite un sujet très pertinent et actuel : la détection du phishing par e-mail à l’aide de techniques de Machine Learning. Le choix du sujet est intéressant, car le hameçonnage représente une menace réelle en cybersécurité, et l’idée de proposer une application permettant d’analyser des fichiers .eml est cohérente avec un besoin concret.
%
% Le rapport contient plusieurs points positifs. La page de garde est claire, les remerciements sont bien rédigés, le résumé et l’Abstract sont présents, la table des matières est propre, et les listes des figures et des tableaux sont incluses. L’introduction générale présente correctement le contexte, la motivation, le problème traité et les objectifs du projet.
%
% La partie contexte est intéressante. Vous expliquez les différentes formes de phishing : phishing de masse, spear-phishing, whaling et techniques modernes d’évasion comme le spoofing, les homoglyphes, le typosquatting et les raccourcisseurs d’URL. Cette partie montre une bonne compréhension du domaine de la cybersécurité.
%
% La partie spécification est claire. Les exigences fonctionnelles sont bien identifiées : gestion des utilisateurs, upload sécurisé de fichiers .eml, extraction en mémoire, extraction des liens, prédiction par IA et génération d’un rapport de risque. Les exigences non fonctionnelles sont également pertinentes, notamment la confidentialité, la politique de Zero Retention, la protection CSRF et le temps de réponse attendu.
%
% La partie Machine Learning constitue l’un des points forts du rapport. Vous expliquez le pipeline utilisé : extraction de caractéristiques structurelles, détection des URLs, mots-clés liés au phishing, vectorisation TF-IDF, entraînement d’un modèle Random Forest, sérialisation avec joblib et intégration dans Django. Les métriques annoncées sont solides : accuracy de 97 %, précision phishing de 94 %, recall phishing de 97 % et ROC-AUC de 0,9928. Ces résultats montrent un vrai effort d’évaluation du modèle.
%
% La partie réalisation est également intéressante. Les choix techniques sont cohérents : Python, Django, PostgreSQL, scikit-learn, pandas, joblib, HTML/CSS et Git/GitHub. Les captures d’écran de l’interface d’upload et du tableau de bord des résultats montrent un travail concret et une interface claire. Les annexes avec les extraits de code sont un très bon point, car elles permettent de vérifier la logique d’extraction, d’entraînement, de prédiction, de parsing .eml, de sécurité et de gestion des rapports.
%
% La bibliographie est satisfaisante. Elle contient plusieurs références scientifiques pertinentes sur la détection du phishing, le Machine Learning, Random Forest, scikit-learn, les datasets de phishing et les rapports de cybersécurité. Elle répond donc mieux aux attentes académiques que beaucoup de bibliographies purement techniques.
%
% Cependant, certaines limites doivent être signalées. D’abord, même si le document compte 51 pages au total, une grande partie correspond aux annexes de code. Le contenu principal du rapport reste relativement court par rapport au minimum demandé hors annexes. Il faudrait enrichir les chapitres principaux, notamment l’état de l’art, la conception, la réalisation et les tests.
%
% La partie conception reste un peu limitée. Le rapport contient un diagramme de cas d’utilisation et un diagramme de classes, mais il manque des diagrammes importants pour mieux comprendre le fonctionnement dynamique du système : diagramme de séquence pour le processus d’upload/analyse, diagramme d’activité du pipeline de détection, ou encore schéma d’architecture détaillé montrant les échanges entre l’utilisateur, Django, PostgreSQL, le parser .eml et le modèle ML.
%
% La partie tests doit aussi être approfondie. Les métriques Machine Learning sont présentes et intéressantes, mais il manque un vrai tableau de tests fonctionnels de l’application web. Il faudrait détailler les scénarios testés : upload d’un fichier .eml valide, fichier invalide, fichier trop volumineux, e-mail sans corps, e-mail sans URL, utilisateur anonyme, utilisateur authentifié, consultation d’un rapport privé, suppression d’un rapport, tentative d’accès non autorisé et gestion des erreurs. Il serait aussi utile d’ajouter une matrice de confusion et un rapport de classification complet pour mieux interpréter les résultats du modèle.
%
% La partie VirusTotal est pertinente, mais elle est présentée comme une perspective d’intégration future. Il faut donc bien distinguer ce qui est réellement implémenté dans la version actuelle et ce qui reste prévu pour une version ultérieure. Cette clarification est importante pour éviter de donner l’impression que l’API VirusTotal est déjà intégrée.
%
% Quelques éléments techniques méritent aussi d’être mieux expliqués. Par exemple, le choix de Random Forest pourrait être comparé avec d’autres modèles comme Logistic Regression, SVM, Naive Bayes ou XGBoost. Le dataset Kaggle utilisé devrait être décrit plus précisément : nombre d’e-mails phishing/safe, équilibre des classes, nettoyage effectué, exemples de données, limites du dataset et risques de généralisation.
%
% Il existe enfin quelques problèmes de forme mineurs : certaines coupures de mots sont visibles dans le PDF, quelques formulations sont très techniques et pourraient être mieux vulgarisées, et certains titres pourraient être harmonisés. Une relecture finale permettrait d’améliorer la fluidité et la présentation académique.
%
% En conclusion, votre rapport présente un projet intéressant, techniquement solide et bien orienté cybersécurité. La partie Machine Learning est sérieuse, l’application Django semble fonctionnelle, les annexes de code renforcent la crédibilité du travail, et la bibliographie est correcte. Les principales améliorations concernent surtout l’enrichissement du contenu principal hors annexes, l’ajout de tests fonctionnels détaillés, l’amélioration de la conception UML et la clarification entre fonctionnalités réalisées et perspectives futures.
%
% Note finale proposée : 15 / 20
%
% Cordialement,

% ============================================================
%  rapport pfa – fakefinder
%  emsi casablanca – 3iir s2
% ============================================================
\documentclass[12pt,a4paper]{report}

% ── encodage & langue ────────────────────────────────────────
\usepackage[utf8]{inputenc}
\usepackage[t1]{fontenc}
\usepackage[french]{babel}

% ── typographie stricte (times new roman & interligne 1.5) ───
\usepackage{mathptmx} 
\usepackage{setspace}
\onehalfspacing       

% ── mise en page (marges 2.5cm) ──────────────────────────────
\usepackage[top=2.5cm, bottom=2.5cm, left=2.5cm, right=2.5cm]{geometry}

% ── graphiques, tableaux & mathématiques ─────────────────────
\usepackage{graphicx}
\usepackage{xcolor}
\usepackage{float}
\usepackage{booktabs}
\usepackage{array}
\usepackage{tabularx}
\usepackage{amsmath}
\usepackage{amssymb}

% ── code source & divers ─────────────────────────────────────
\usepackage{listings}
\usepackage{enumitem}
\usepackage{titlesec}
\usepackage[hidelinks]{hyperref}

% ── couleurs ─────────────────────────────────────────────────
\definecolor{codegray}{rgb}{0.95,0.95,0.95}
\definecolor{codeblue}{rgb}{0.18,0.27,0.60}
\definecolor{codegreen}{rgb}{0,0.6,0}
\definecolor{codepurple}{rgb}{0.58,0,0.82}
\definecolor{backcolour}{rgb}{0.95,0.95,0.92}

\lstset{
    backgroundcolor=\color{backcolour},
    commentstyle=\color{codegreen},
    keywordstyle=\color{magenta},
    numberstyle=\tiny\color{codegray},
    stringstyle=\color{codepurple},
    basicstyle=\ttfamily\footnotesize,
    breakatwhitespace=false,
    breaklines=true,
    captionpos=b,
    keepspaces=true,
    numbers=left,
    numbersep=5pt,
    showspaces=false,
    showstringspaces=false,
    showtabs=false,
    tabsize=2
}

% format des titres de chapitres (taille 14 selon cdc emsi)
\titleformat{\chapter}[display]
  {\normalfont\fontsize{14}{16}\bfseries}
  {\chaptertitlename\ \thechapter}{10pt}
  {\fontsize{14}{16}\bfseries}

% tolérance pour éviter les mots coupés de façon disgracieuse
\sloppy 

% ============================================================
\begin{document}

% ── numérotation romaine pour les pages préliminaires ────────
\pagenumbering{roman}
\thispagestyle{empty}

% ============================================================
%  1. page de garde
% ============================================================
\begin{titlepage}
  \centering

  \includegraphics[width=0.4\textwidth]{emsi_logo.png}\\[0.5cm]
  \textbf{\large école marocaine des sciences de l'ingénieur}\\
  \textbf{\large emsi casablanca}\\[0.3cm]
  \rule{\linewidth}{0.8pt}\\[0.3cm]
  \textit{filière : ingénierie informatique et réseaux – 3iir}\\
  \textit{année universitaire : 2025/2026}\\[1.5cm]

  \rule{\linewidth}{1.5pt}\\[0.4cm]
  {\fontsize{18}{22}\bfseries
    rapport de projet de fin d'année\\[0.3cm]
    \textcolor{codeblue}{\fontsize{20}{24}\selectfont fakefinder}\\[0.2cm]
    \fontsize{14}{18}\normalfont
    système intelligent de détection du hameçonnage par e-mail
  }\\[0.4cm]
  \rule{\linewidth}{1.5pt}\\[2cm]

  \begin{minipage}{0.48\textwidth}
    \raggedright
    \textbf{réalisé par :}\\[0.2cm]
    yahya khalifi\\
    elias chakroun\\[0.8cm]
  \end{minipage}
  \hfill
  \begin{minipage}{0.48\textwidth}
    \raggedleft
    \textbf{tuteur académique :}\\[0.2cm]
    mme soukayna mayara\\[0.8cm]
    \textbf{date de soutenance :}\\[0.2cm]
    18 mai 2026
  \end{minipage}

  \vfill
  \rule{\linewidth}{0.4pt}\\[0.2cm]
  {\small membre de honoris united universities}
\end{titlepage}

% ============================================================
%  2. remerciements
% ============================================================
\chapter*{remerciements}
\addcontentsline{toc}{chapter}{remerciements}

la réalisation de ce projet de fin d'année a été une expérience profondément enrichissante, tant sur le plan technique que professionnel. nous tenons à exprimer notre profonde gratitude envers toutes les personnes qui ont contribué, de près ou de loin, à l'aboutissement de ce travail de recherche et de développement.

nos remerciements s'adressent en tout premier lieu à notre encadrante pédagogique, \textbf{madame soukayna mayara}. son expertise technique, sa rigueur scientifique, ses conseils avisés et sa grande disponibilité ont été des atouts majeurs tout au long de ce projet. ses directives nous ont permis de structurer notre travail de manière méthodique, de surmonter les obstacles techniques et de maintenir un haut niveau d'exigence académique.

nous souhaitons également témoigner notre reconnaissance au corps professoral et à l'administration de \textbf{l'école marocaine des sciences de l'ingénieur (emsi)} de casablanca. la qualité de l'enseignement dispensé, la pertinence des modules abordés durant notre cursus en ingénierie informatique et réseaux, ainsi que les infrastructures mises à notre disposition ont constitué le socle indispensable à la réalisation de ce projet complexe alliant cybersécurité et intelligence artificielle.

enfin, nous adressons nos pensées les plus chaleureuses à nos familles, qui n'ont cessé de nous encourager, et à nos amis et camarades de promotion pour leur soutien moral indéfectible, leur patience, et les moments de partage et d'entraide qui ont jalonné ces années d'études.

% ============================================================
%  3. résumé & abstract
% ============================================================
\chapter*{résumé}
\addcontentsline{toc}{chapter}{résumé}

le projet \textbf{fakefinder} est une solution de sécurité logicielle novatrice, développée autour du framework web django, conçue spécifiquement pour lutter contre l'une des cybermenaces les plus pernicieuses et répandues de notre décennie : le hameçonnage (phishing) par courrier électronique. face à la sophistication grandissante des attaques d'ingénierie sociale, qui parviennent régulièrement à tromper la vigilance des filtres antispam traditionnels basés sur de simples listes noires, notre système propose une approche proactive fondée sur l'intelligence artificielle.

en s'appuyant sur des modèles de machine learning avancés, et plus particulièrement sur l'algorithme des forêts aléatoires (random forest) couplé à une vectorisation sémantique tf-idf (term frequency-inverse document frequency), ce système effectue une analyse sémantique, lexicale et structurelle approfondie des en-têtes et du corps des e-mails. l'application offre aux utilisateurs finaux une interface web intuitive et ergonomique permettant de vérifier la légitimité des fichiers bruts de messagerie (format \texttt{.eml}) de manière quasi instantanée. 

l'un des piliers majeurs de ce projet réside dans son architecture de confidentialité par conception (privacy by design). fakefinder applique une politique stricte de "zero retention", garantissant que les données sensibles contenues dans les e-mails soumis sont traitées intégralement en mémoire vive (ram) et détruites immédiatement après l'inférence du modèle, sans jamais être écrites sur un support de stockage persistant.

ce rapport détaille de manière exhaustive le processus de conception architecturale, la phase cruciale d'ingénierie des caractéristiques (feature engineering), l'implémentation technique du backend et du pipeline de données, ainsi que l'évaluation fonctionnelle et statistique du système. il explore par ailleurs les évolutions futures de la plateforme, avec une mise en lumière particulière sur l'intégration prévue d'interfaces de programmation (api) de threat intelligence telles que virustotal, visant à offrir une vérification dynamique et exhaustive de la réputation des url extraites.

\textbf{mots-clés :} hameçonnage, cybersécurité, intelligence artificielle, machine learning, random forest, tf-idf, django, python, zero retention, ingénierie sociale.

\vspace{1cm}

\chapter*{abstract}
\addcontentsline{toc}{chapter}{abstract}

the \textbf{fakefinder} project is an innovative software security solution built around the django web framework, specifically designed to combat one of the most pervasive and damaging cyber threats of our decade: email phishing. in the face of increasingly sophisticated social engineering attacks that frequently bypass traditional, blacklist-based spam filters, our system proposes a proactive approach founded on artificial intelligence.

by leveraging advanced machine learning models—specifically the random forest algorithm coupled with tf-idf (term frequency-inverse document frequency) semantic vectorization—this system performs an in-depth semantic, lexical, and structural analysis of both email headers and body content. the application provides end-users with an intuitive and ergonomic web interface to verify the legitimacy of raw email files (\texttt{.eml} format) almost instantaneously.

a major pillar of this project lies in its "privacy by design" architecture. fakefinder enforces a strict "zero retention" policy, ensuring that the sensitive data contained in the submitted emails is processed entirely in random access memory (ram) and destroyed immediately after model inference, without ever being written to persistent storage.

this report exhaustively details the architectural design process, the crucial feature engineering phase, the technical implementation of the backend and data pipeline, as well as the functional and statistical evaluation of the system. furthermore, it explores future developments for the platform, with a particular focus on the planned integration of threat intelligence application programming interfaces (apis) such as virustotal, aimed at providing dynamic and comprehensive reputation verification for extracted urls.

\textbf{keywords:} phishing, cybersecurity, artificial intelligence, machine learning, random forest, tf-idf, django, python, zero retention, social engineering.

% ============================================================
%  4. tables et listes
% ============================================================
\tableofcontents
\newpage
\listoffigures
\addcontentsline{toc}{chapter}{table des figures}
\newpage
\listoftables
\addcontentsline{toc}{chapter}{liste des tableaux}

% ============================================================
%  introduction générale
% ============================================================
\cleardoublepage
\pagenumbering{arabic}
\setcounter{page}{1}

\chapter*{introduction générale}
\addcontentsline{toc}{chapter}{introduction générale}

\section*{contexte global et motivation}
dans le paysage numérique contemporain, la messagerie électronique demeure, malgré l'émergence des messageries instantanées, la pierre angulaire de la communication professionnelle et personnelle à l'échelle mondiale. avec plus de 300 milliards d'e-mails échangés quotidiennement, ce canal de communication universel est naturellement devenu le terrain de chasse privilégié et le principal vecteur d'attaque exploité par les cybercriminels.

le hameçonnage (ou phishing), une forme pernicieuse et psychologique d'ingénierie sociale, consiste à usurper l'identité numérique d'une entité de confiance (institution bancaire, administration gouvernementale, fournisseur de services cloud, ou même un collègue de travail) dans le but de tromper une cible. l'objectif final est souvent double : soutirer des informations sensibles (identifiants de connexion, données bancaires, secrets industriels) ou inciter la victime à télécharger et exécuter, à son insu, des charges utiles malveillantes (ransomwares, chevaux de troie, spywares).

au cours de la dernière décennie, le hameçonnage a connu une mutation radicale. il a évolué de simples arnaques de masse, souvent truffées de fautes d'orthographe et grossières dans leur conception, vers des campagnes hautement ciblées et d'une sophistication redoutable, connues sous le nom de \textit{spear-phishing} ou de \textit{whaling}. les attaquants exploitent désormais l'intelligence artificielle générative pour rédiger des messages sans faille et utilisent des techniques d'obfuscation structurelles avancées, contournant ainsi allègrement les passerelles de sécurité traditionnelles (secure email gateways) déployées par les entreprises.

la motivation première derrière le projet \textbf{fakefinder} découle directement de ce constat alarmant, mis en exergue par des rapports mondiaux tels que le \textit{data breach investigations report} de verizon, qui place systématiquement le phishing dans le top 3 des causes initiales de compromission de données. il existe aujourd'hui un besoin critique pour un outil accessible au grand public et aux professionnels, indépendant de leur fournisseur de messagerie, et suffisamment puissant pour permettre de vérifier un e-mail suspect avant toute interaction compromettante. fakefinder se positionne comme un système de "deuxième avis" diagnostique, effectuant une analyse approfondie et algorithmique des fichiers de messagerie bruts.

\section*{énoncé de la problématique}
la sécurité défensive traditionnelle des e-mails s'appuie historiquement sur des approches réactives. les filtres antispam fonctionnent majoritairement via des listes noires (blacklists) recensant les adresses ip et les noms de domaine réputés malveillants, ou via la comparaison de signatures de malwares connus. or, ces méthodes sont structurellement inefficaces contre les attaques de type \textit{zero-day} (nouvelles menaces fraîchement déployées et non encore répertoriées par la communauté cyber). de plus, lorsque des attaquants compromettent des comptes e-mails parfaitement légitimes pour envoyer leurs messages frauduleux, la vérification de la réputation de l'expéditeur devient obsolète.

par ailleurs, la grande majorité des utilisateurs finaux, qui constituent le maillon le plus vulnérable de la chaîne de sécurité, ne disposent ni des compétences techniques requises ni du temps matériel nécessaire pour inspecter manuellement les en-têtes de routage cryptiques d'un e-mail (tels que les champs \texttt{received}, \texttt{authentication-results}, ou \texttt{return-path}) à la recherche d'indicateurs d'usurpation.

la problématique centrale de ce projet peut donc se formuler ainsi : \textbf{comment concevoir un système automatisé, performant et respectueux de la confidentialité, capable de détecter les e-mails de phishing de nouvelle génération en analysant les subtilités sémantiques et structurelles du message, sans s'appuyer exclusivement sur des bases de données de menaces connues ?}

\section*{objectifs et portée du projet}
pour répondre à cette problématique, ce projet de fin d'année (pfa) vise à développer une solution hybride baptisée \textbf{fakefinder}, à la croisée du développement logiciel web et de l'intelligence artificielle. les objectifs stratégiques et techniques se déclinent comme suit :

\begin{itemize}
    \item \textbf{ingénierie de traitement des données brutes :} développer un moteur d'analyse syntaxique (parser) hautement performant en python, capable de disséquer la structure complexe (mime) d'un fichier \texttt{.eml}. ce moteur doit extraire le corps du texte, identifier les liens hyperboliques, et isoler les métadonnées, quels que soient les encodages ou les méthodes d'obfuscation utilisés par l'attaquant.
    \item \textbf{modélisation par apprentissage automatique (machine learning) :} concevoir, entraîner et déployer un classificateur robuste. le choix s'est porté sur l'algorithme des \textit{forêts aléatoires (random forest)}, optimisé pour la détection sémantique de textes frauduleux via la méthode mathématique de vectorisation \textit{tf-idf}.
    \item \textbf{garantie absolue de confidentialité (zero retention) :} concevoir une architecture garantissant que le système traite les données de manière strictement éphémère. aucun fichier soumis, ni le texte de son contenu, ne doit être persisté sur les disques ou la base de données du serveur, assurant ainsi une conformité totale avec les réglementations sur la protection des données (ex: rgpd).
    \item \textbf{développement d'une plateforme web full-stack :} créer une interface web ergonomique, responsive et sécurisée à l'aide du framework django, simplifiant l'expérience utilisateur grâce à des fonctionnalités modernes (comme le glisser-déposer de fichiers).
    \item \textbf{scalabilité et évolutivité architecturale :} préparer le code source et la structure logicielle à des intégrations futures majeures, notamment le requêtage d'api de \textit{threat intelligence} (comme virustotal) pour croiser l'analyse sémantique interne avec des renseignements mondiaux sur les cybermenaces.
\end{itemize}

\section*{structure du rapport}
ce rapport est structuré en six chapitres distincts, reflétant le cycle de vie complet du développement du projet :
\begin{itemize}
    \item le \textbf{chapitre 1} présente l'établissement académique d'accueil (emsi) et plonge dans le contexte technologique de la menace du phishing, en détaillant l'anatomie d'une attaque.
    \item le \textbf{chapitre 2} aborde la phase de spécification du projet, détaillant les exigences fonctionnelles, non fonctionnelles, et la gestion du projet via la méthodologie agile.
    \item le \textbf{chapitre 3} est consacré à l'analyse et à la conception du système, incluant les modélisations uml et l'architecture logicielle retenue.
    \item le \textbf{chapitre 4} constitue le cœur technique du rapport. il détaille la réalisation, le développement du backend, le pipeline de traitement des données, ainsi que la théorie et l'implémentation du modèle d'intelligence artificielle.
    \item le \textbf{chapitre 5} explore les perspectives d'évolution de la plateforme, avec un focus sur la mécanique d'intégration future de services de renseignements sur les menaces (virustotal).
    \item enfin, le \textbf{chapitre 6} présente les résultats des tests fonctionnels et l'évaluation statistique du modèle de machine learning, clôturant le rapport par un bilan global du projet.
\end{itemize}


% ============================================================
%  chapitre 1 : contexte et école
% ============================================================
\chapter{Présentation de l'Organisme et Contexte de la Menace}

\section{présentation de l'emsi}
l'école marocaine des sciences de l'ingénieur (emsi) est un établissement privé d'enseignement supérieur fondé en 1986. reconnue par l'état marocain, l'emsi s'est progressivement imposée au fil des décennies comme la plus grande école d'ingénieurs privée du royaume. sa mission principale est de former des ingénieurs d'état et des cadres pluridisciplinaires, capables d'accompagner le développement technologique, industriel et économique national et international.

l'établissement se distingue par son approche pédagogique pragmatique et orientée vers les besoins réels du marché du travail. membre fondateur du réseau panafricain \textit{honoris united universities}, l'école met un point d'honneur à offrir une formation qui allie théorie académique rigoureuse et pratique intensive par le biais de laboratoires, de stages industriels et de projets de recherche.

l'emsi structure ses formations autour de plusieurs filières d'ingénierie d'excellence. parmi celles-ci, la filière \textbf{ingénierie informatique et réseaux (iir)} vise à doter les futurs ingénieurs de compétences pointues et transversales : conception et développement de systèmes d'information complexes, administration de bases de données, gestion d'infrastructures réseaux, intelligence artificielle, et cybersécurité.

c'est dans ce cadre exigeant et professionnalisant, au cours de la troisième année de la filière iir (3iir), qu'intervient ce projet de fin d'année (pfa). ce projet constitue une mise en pratique concrète et une synthèse des modules théoriques de sécurité des systèmes d'information, de génie logiciel et d'analyse de données étudiés tout au long du cursus universitaire.

\begin{figure}[h]
    \centering
    \includegraphics[width=0.7\textwidth]{organigramme_emsi.png}
    \caption{structure organisationnelle et représentation des pôles d'excellence de l'emsi}
\end{figure}

\section{anatomie et taxonomie de la menace du phishing}
pour concevoir un système de défense et de détection qui soit véritablement efficace, il est une condition \textit{sine qua non} : comprendre en profondeur la nature, l'évolution et les mécanismes internes de la menace que l'on cherche à contrer. le phishing, dans son acception moderne, n'est plus une attaque monolithique ; il se décline aujourd'hui en plusieurs vecteurs hautement spécialisés, exploitant à la fois des failles techniques et des biais cognitifs humains.

\subsection{typologie des attaques}
les acteurs malveillants adaptent leurs stratégies en fonction de la cible et des ressources à leur disposition. on distingue aujourd'hui trois grandes catégories de hameçonnage :

\begin{itemize}
    \item \textbf{le phishing de masse (deceptive phishing) :} il s'agit de la forme la plus ancienne et la plus classique. l'attaquant envoie un volume massif d'e-mails génériques (souvent des millions par campagne) en usurpant l'identité d'organisations de grande envergure et grand public (opérateurs télécoms, plateformes de streaming comme netflix, banques de détail, services postaux ou de livraison comme dhl). le modèle repose sur la loi des grands nombres : l'objectif est purement statistique. même avec un taux de conversion (clic) inférieur à 0,1\%, la campagne s'avère rentable. l'urgence est souvent le levier principal : "votre colis n'a pas pu être livré", "votre abonnement est suspendu".
    
    \item \textbf{le spear-phishing (harponnage) :} contrairement au phishing de masse, le spear-phishing est chirurgical et hautement qualitatif. l'attaquant cible un individu spécifique, un département particulier (ex: les ressources humaines), ou un groupe restreint au sein d'une organisation. cette attaque est invariablement précédée par une minutieuse phase de reconnaissance et de collecte de renseignements en sources ouvertes (osint - open source intelligence). en fouillant les réseaux sociaux professionnels (linkedin), les publications de l'entreprise ou les fuites de données antérieures, l'attaquant construit un profil de la victime. l'e-mail frauduleux contient ainsi des éléments de contexte ultra-personnalisés (nom du manager, référence à un projet en cours, vocabulaire interne), rendant la supercherie indétectable pour un œil non averti.
    
    \item \textbf{le whaling (chasse à la baleine) :} le whaling est une sous-catégorie très spécifique du spear-phishing, ciblant exclusivement les individus possédant un fort pouvoir décisionnel ou financier au sein de l'organisation : membres du comité de direction, pdg (ceo), directeurs financiers (cfo). le préjudice visé est très souvent financier (connue sous le nom de "fraude au président" ou bec - business email compromise). l'attaquant, se faisant passer pour le dirigeant légitime via une adresse usurpée, ordonne par exemple un virement urgent, confidentiel et exceptionnel à un employé du service comptabilité en vue d'une "acquisition secrète".
\end{itemize}

\subsection{anatomie technique d'un e-mail et techniques d'évasion}
un courrier électronique n'est pas qu'un simple texte ; c'est un ensemble structuré de données codées selon la norme mime (multipurpose internet mail extensions) et transférées via le protocole smtp (simple mail transfer protocol). pour contourner les filtres de sécurité, les attaquants manipulent ces structures sous-jacentes.

\textbf{1. l'usurpation d'identité et les en-têtes smtp (spoofing) :}
l'architecture originelle du protocole smtp, datant des années 1980, ne prévoyait aucun mécanisme robuste d'authentification de l'expéditeur. il est techniquement trivial d'écrire un script qui envoie un e-mail où le champ d'en-tête \texttt{from:} affiche \texttt{admin@paypal.com}, bien que l'e-mail provienne d'un serveur malveillant en russie. pour contrer cela, des protocoles ont été créés (spf, dkim, dmarc), mais de nombreuses entreprises ne les configurent pas correctement, laissant la porte ouverte au \textit{direct spoofing}.

\textbf{2. homoglyphes et typosquatting (attaques sémantiques) :}
lorsque le domaine cible est protégé par dmarc, les attaquants utilisent le typosquatting : ils enregistrent un nom de domaine visuellement très proche du domaine légitime. par exemple :
\begin{itemize}
    \item inversion ou oubli de lettres : \texttt{secuite-societegenerale.fr} au lieu de \texttt{securite-societegenerale.fr}.
    \item substitution visuelle : \texttt{paypa1.com} (avec un "un" au lieu d'un "l").
    \item homoglyphes (attaque idn) : remplacement de caractères latins par des caractères d'autres alphabets (ex: le "a" latin remplacé par le "a" cyrillique u+0430). visuellement identiques pour l'humain, ils sont mathématiquement différents pour la machine.
\end{itemize}

\textbf{3. obfuscation des liens (payload delivery) :}
le lien contenu dans l'e-mail est l'arme du crime. pour éviter que les antivirus n'analysent la destination et ne bloquent le message, les attaquants utilisent :
\begin{itemize}
    \item des \textbf{raccourcisseurs d'url} légitimes (\texttt{bit.ly}, \texttt{tinyurl.com}) pour masquer la destination finale.
    \item des \textbf{redirections ouvertes (open redirects) :} exploitation d'une faille sur un site légitime très bien réputé (ex: \texttt{https://google.com/url?q=http://site-hacker.com}), ce qui trompe les filtres qui ne vérifient que le premier domaine.
    \item la \textbf{stéganographie html :} cacher du code malveillant dans des balises html invisibles, ou séparer un lien malveillant avec des commentaires html pour tromper l'analyse syntaxique des passerelles de sécurité, tout en s'affichant normalement dans le navigateur de la victime.
\end{itemize}

\section{la réponse de fakefinder}
face à ces techniques qui fusionnent une ingénierie sociale raffinée (pour tromper l'humain) et des astuces techniques complexes (pour tromper les filtres basés sur des règles fixes), une approche statique est vouée à l'échec. la réponse défensive proposée par fakefinder consiste à ne plus chercher des "signatures" exactes d'attaques passées, mais à utiliser le \textbf{machine learning} pour identifier les "modèles" comportementaux, le champ lexical de l'urgence, et les anomalies structurelles globales du message. l'intelligence artificielle est ainsi employée pour déceler les corrélations mathématiques invisibles à l'analyse humaine standard.


% ============================================================
%  chapitre 2 : spécification et gestion du projet
% ============================================================
\chapter{Spécification et Gestion du Projet}

avant toute ligne de code, le développement d'un logiciel robuste nécessite une définition rigoureuse de ce que le système doit accomplir, de la manière dont il doit se comporter, et des contraintes dans lesquelles il doit s'inscrire. ce chapitre détaille le cahier des charges de fakefinder et la méthodologie employée pour mener à bien sa réalisation.

\section{ingénierie des exigences fonctionnelles (fr)}
les exigences fonctionnelles décrivent les actions spécifiques que le système doit exécuter pour apporter de la valeur à l'utilisateur. le tableau ci-dessous formalise les fonctionnalités maîtresses attendues de la plateforme.

\begin{table}[h]
\centering
\begin{tabular}{|l|p{12cm}|}
\hline
\textbf{id} & \textbf{description détaillée de l'exigence fonctionnelle} \\ \hline
fr-1 & \textbf{gestion des profils et authentification :} le système doit offrir un module complet permettant aux utilisateurs de créer un compte, de s'authentifier de manière sécurisée (hachage des mots de passe) et de gérer leur session. \\ \hline
fr-2 & \textbf{espace de téléversement (upload) sécurisé :} l'application doit fournir une interface web ergonomique (drag-and-drop) pour téléverser des fichiers au format standard de messagerie rfc 822 (\texttt{.eml}). une limite stricte de taille de fichier (10 mo) doit être appliquée en amont pour prévenir les attaques de type déni de service (dos). \\ \hline
fr-3 & \textbf{moteur d'extraction et de parsing in-memory :} le backend du système doit être capable de réceptionner le fichier \texttt{.eml} directement en mémoire vive, d'en analyser la structure mime complexe, d'extraire les en-têtes pertinents (expéditeur, domaine, sujet) et de décoder le corps du message (qu'il soit en \texttt{text/plain} ou en \texttt{text/html}) sans nécessiter de sauvegarde temporaire sur le disque du serveur. \\ \hline
fr-4 & \textbf{filtrage et extraction des liens (url) :} lors de l'analyse du corps du texte, le système doit exécuter des expressions régulières (regex) avancées pour identifier, isoler et assainir la totalité des liens hyperboliques contenus dans le message. \\ \hline
fr-5 & \textbf{inférence par intelligence artificielle :} les données textuelles extraites doivent être nettoyées, transformées en vecteurs mathématiques (tf-idf), puis injectées dans le modèle de machine learning pré-entraîné pour obtenir en retour une probabilité calculée de fraude (classification binaire). \\ \hline
fr-6 & \textbf{génération et restitution du rapport d'analyse :} le système doit consolider toutes les informations recueillies et la prédiction de l'ia pour générer un rapport dynamique. ce rapport doit présenter : un score global sur 100, un code couleur d'alerte intuitif (vert/orange/rouge correspondant à faible/moyen/élevé), le domaine réel de l'expéditeur, et la liste exhaustive des url suspectes identifiées. \\ \hline
fr-7 & \textbf{historisation (utilisateurs connectés) :} si l'utilisateur est authentifié lors de la soumission, le système doit enregistrer les métadonnées du rapport d'analyse dans une base de données, permettant à l'utilisateur de consulter l'historique de ses contrôles via un tableau de bord privé. \\ \hline
\end{tabular}
\caption{tableau exhaustif des exigences fonctionnelles (functional requirements)}
\end{table}

\section{exigences non-fonctionnelles (nfr)}
les exigences non-fonctionnelles ne décrivent pas *ce que* fait le système, mais *comment* il le fait. elles définissent les attributs de qualité, cruciaux pour un outil de cybersécurité.

\subsection{1. contraintes de sécurité et de confidentialité : architecture "zero retention"}
l'exigence non-fonctionnelle la plus critique du projet concerne la protection des données. les e-mails professionnels ou personnels soumis par les utilisateurs peuvent contenir des informations hautement confidentielles (mots de passe, données financières, secrets d'entreprise, correspondances privées). 
il est donc hors de question que l'application fakefinder devienne elle-même un risque pour la vie privée. 

pour y répondre, nous avons conçu l'architecture autour d'un principe absolu de \textbf{zero retention} :
\begin{itemize}
    \item \textbf{traitement éphémère en ram :} dès sa réception via la requête http post, le fichier \texttt{.eml} est manipulé exclusivement sous forme d'objet mémoire (`inmemoryuploadedfile` sous django).
    \item \textbf{aucune persistance disque :} le contenu brut de l'e-mail, ainsi que le texte extrait du corps du message, ne sont \textbf{jamais} écrits sur le disque dur du serveur, et ne sont \textbf{jamais} insérés dans une table de la base de données.
    \item \textbf{nettoyage actif :} dès que le moteur d'intelligence artificielle a généré son score, la variable contenant le texte de l'e-mail est libérée, et le fichier temporaire en mémoire est détruit par un \textit{middleware} django de sécurité personnalisé fonctionnant en post-traitement de la requête.
    \item \textbf{stockage limité aux métadonnées :} la base de données postgresql enregistre uniquement les résultats de l'analyse (le score, la couleur, la date) et des indicateurs de menace dépersonnalisés (les url extraites et le nom de domaine de l'expéditeur), rendant impossible la reconstruction du message originel en cas de compromission de la base de données.
\end{itemize}

\subsection{2. sécurité applicative web}
le développement reposant sur django bénéficie de protections natives robustes contre les vulnérabilités majeures du web (top 10 owasp) :
\begin{itemize}
    \item protection contre les attaques \textbf{csrf (cross-site request forgery)} obligatoire sur tous les formulaires modifiant l'état.
    \item protection \textbf{xss (cross-site scripting)} grâce à l'échappement automatique (auto-escaping) du moteur de gabarits (templates) django, particulièrement important lors de l'affichage des url extraites des e-mails malveillants.
    \item prévention des injections \textbf{sql} grâce à l'utilisation systématique de l'orm (object-relational mapping) de django, évitant les requêtes en base de données concaténées manuellement.
\end{itemize}

\subsection{3. performances, disponibilité et temps de réponse}
un outil de sécurité n'est utilisé que s'il est fluide. l'expérience utilisateur (ux) exige que le processus complet de vérification n'entraîne pas de friction majeure. l'exigence de performance fixe un seuil maximal : l'analyse complète d'un e-mail de taille standard (inférieure à 1 mo) — processus incluant l'envoi http, le parsing complexe de l'arborescence mime, la tokenization nlp, l'inférence du modèle random forest, et le rendu html du résultat — doit s'exécuter de manière synchrone en **moins de 5 secondes**.
pour atteindre cet objectif, le modèle machine learning (fichier \texttt{.joblib}) doit être chargé en mémoire de manière paresseuse (lazy loading) au démarrage du serveur, évitant ainsi le coût d'une lecture disque à chaque requête utilisateur.

\section{planification et gestion de projet (méthodologie agile)}
pour assurer un développement itératif, incrémental et maîtrisé sur un calendrier restreint de 5 semaines, le projet a été piloté en s'inspirant fortement de la méthodologie **agile (framework scrum)**. cette approche, bien qu'adaptée pour une équipe réduite (binôme), a permis de définir des cycles de développement courts appelés sprints, chacun délivrant un incrément fonctionnel du produit.

\begin{table}[h]
\centering
\caption{planification détaillée des sprints de réalisation (avril - mai 2026)}
\begin{tabularx}{\textwidth}{|c|x|l|}
\hline
\textbf{semaine} & \textbf{périmètre du sprint et tâches de développement} & \textbf{livrable du sprint} \\ \hline
\textbf{s1} & \textbf{phase de lancement \& conception :} revue de littérature sur les modèles ml en cybersécurité, identification et collecte du dataset adéquat sur la plateforme kaggle, spécification des exigences et création des diagrammes uml initiaux. & dossier de conception \\ \hline
\textbf{s2} & \textbf{développement du core backend :} initialisation du projet django, modélisation de la base de données postgresql (entité scanreport). développement intensif du script utilitaire de parsing \texttt{.eml} en mémoire vive. & api backend / parseur \\ \hline
\textbf{s3} & \textbf{ingénierie ia \& machine learning :} scripting du pipeline de préparation des données textuelles. nettoyage du dataset, configuration de la vectorisation tf-idf. choix, entraînement et calibration de l'algorithme random forest. sauvegarde du modèle. & modèle sérialisé (\texttt{.joblib}) \\ \hline
\textbf{s4} & \textbf{intégration \& développement frontend :} chargement du modèle ia dans les vues django. conception des maquettes de l'interface utilisateur. intégration html/css (responsive design) pour les pages d'upload et les tableaux de bord de résultats. & version beta fonctionnelle \\ \hline
\textbf{s5} & \textbf{validation, optimisation \& clôture :} débogage intensif, exécution de tests fonctionnels et de tests de robustesse (fichiers corrompus). amélioration des performances de chargement. rédaction du rapport académique final. & produit final \& rapport \\ \hline
\end{tabularx}
\end{table}

\subsection{gestion des risques}
durant la phase d'initialisation (sprint 1), une identification des risques inhérents au projet a été réalisée afin d'anticiper les mesures d'atténuation :
\begin{itemize}
    \item \textbf{risque technologique (parsing mime) :} la norme rfc 822 décrivant les e-mails est complexe et très permissive. des e-mails mal formatés pourraient faire crasher le parseur. \textit{atténuation :} utilisation intensive de blocs \texttt{try/except} robustes et de décodeurs d'encodage dynamiques.
    \item \textbf{risque ia (surapprentissage) :} le modèle ml pourrait apprendre par cœur le dataset d'entraînement et échouer sur de nouveaux e-mails. \textit{atténuation :} isolation stricte de 20\% des données (jeu de test) pour évaluer la capacité de généralisation réelle du modèle, et limitation de la profondeur des arbres du random forest.
\end{itemize}


% ============================================================
%  chapitre 3 : analyse et conception
% ============================================================
\chapter{analyse et conception du système}

la phase de conception est la passerelle entre l'expression du besoin (spécifications) et la production du code (développement). elle consiste à définir l'architecture globale, la dynamique des interactions et la structure des données du système, en s'appuyant sur les standards de la modélisation logicielle (uml - unified modeling language).

\section{modélisation de l'architecture applicative}

\subsection{le patron d'architecture mvt (django)}
l'application fakefinder est structurée selon le patron d'architecture **mvt (model-view-template)**, qui est l'implémentation spécifique au framework django du célèbre paradigme mvc (model-view-controller). cette architecture sépare strictement les préoccupations logicielles, garantissant ainsi un code maintenable et évolutif.

\begin{itemize}
    \item \textbf{le modèle (model) :} c'est la couche d'abstraction de la base de données. dans fakefinder, les modèles définissent la structure des tables postgresql sous forme de classes python. c'est le modèle qui dicte que le champ \texttt{risk\_score} est une chaîne de caractères et que \texttt{suspicious\_urls} est un champ json. l'orm (object-relational mapper) de django traduit ensuite ces classes en requêtes sql sécurisées.
    
    \item \textbf{la vue (view) :} correspondant au contrôleur dans le paradigme classique, la vue est le cerveau de l'application. elle réceptionne la requête http entrante (ex: l'upload du fichier \texttt{.eml}). c'est au sein de la vue qu'est orchestrée la logique métier : elle appelle les utilitaires de parsing, invoque le modèle de machine learning pour réaliser l'inférence, demande au modèle d'enregistrer les résultats, et détermine finalement quelle interface renvoyer à l'utilisateur.
    
    \item \textbf{le gabarit (template) :} c'est la couche de présentation. les gabarits sont des fichiers html enrichis d'un langage de balisage spécifique (django template language). la vue passe un "contexte" (des données, comme le score ou la liste des url) au gabarit, qui génère dynamiquement la page html finale envoyée au navigateur web de l'utilisateur.
\end{itemize}

\subsection{schéma d'architecture système et flux de données}
l'architecture de fakefinder repose sur une interaction fluide entre plusieurs composants spécialisés, orchestrés par le framework django :

\begin{enumerate}
    \item \textbf{couche client (frontend) :} l'utilisateur interagit avec une interface web responsive. le fichier \texttt{.eml} est transmis via une requête \texttt{multipart/form-data}.
    \item \textbf{couche applicative (django backend) :}
        \begin{itemize}
            \item \textit{moteur de parsing (parser) :} décompose l'e-mail en mémoire vive, extrait les métadonnées (émetteur, sujet) et le corps du message (texte brut/html).
            \textit{pipeline de feature engineering :} transforme le texte extrait en un vecteur numérique compatible avec le modèle ml.
            \textit{moteur d'inférence (ml model) :} charge le modèle \texttt{.joblib} et calcule la probabilité de phishing.
        \end{itemize}
    \item \textbf{couche de persistance (postgresql) :} stocke uniquement les métadonnées dépersonnalisées et les scores d'analyse, garantissant la politique de zero retention pour le contenu des e-mails.
    \item \textbf{couche de performance (redis) :} fait office de moteur de cache distribué pour les domaines de confiance et les sessions, réduisant drastiquement les accès disque redondants. Un mécanisme de \textit{fallback} transparent a été implémenté pour garantir la continuité du service en cas d'indisponibilité du serveur Redis, basculant automatiquement sur des requêtes directes en base de données.
\end{enumerate}

ce flux assure que les données sensibles ne quittent jamais la mémoire vive du processus applicatif, minimisant ainsi la surface d'attaque et garantissant une confidentialité maximale.

\section{modélisation des comportements et interactions (uml)}

\subsection{diagramme de cas d'utilisation}
le diagramme de cas d'utilisation définit de manière macroscopique les frontières du système et recense les objectifs des acteurs qui interagissent avec lui.

\begin{figure}[h]
    \centering
    \includegraphics[width=0.85\textwidth]{usecase_diagram.png}
    \caption{diagramme de cas d'utilisation illustrant les interactions des acteurs avec le système}
\end{figure}

dans fakefinder, nous identifions deux catégories d'acteurs :
\begin{enumerate}
    \item \textbf{l'utilisateur standard :} il peut utiliser le service de manière anonyme pour soumettre rapidement un fichier \texttt{.eml} et consulter le rapport de dangerosité. s'il choisit de créer un compte et de s'authentifier, l'étendue de ses permissions s'élargit : il a alors la possibilité d'accéder à son tableau de bord personnel, qui liste l'historique complet de ses analyses passées, et d'effacer ses anciens rapports.
    \item \textbf{l'administrateur système :} il s'agit d'un super-utilisateur disposant d'un accès à l'interface d'administration sécurisée de django. ses cas d'utilisation incluent la gestion (création/suppression/bannissement) des comptes utilisateurs, ainsi que la supervision et l'audit de tous les rapports d'analyse générés par le système, quel que soit leur auteur.
\end{enumerate}

\subsection{diagramme de séquence : le flux d'analyse}
le diagramme de séquence est fondamental pour comprendre la chronologie des événements et les communications entre les différents objets du système lors du scénario principal : la soumission d'un e-mail pour analyse.

\begin{figure}[h]
    \centering
    \includegraphics[width=0.85\textwidth]{image_diagramme_sequence.png}
    \caption{diagramme de séquence : parcours temporel de la requête d'analyse}
\end{figure}

**chronologie de l'exécution :**
1. l'utilisateur sélectionne un fichier \texttt{.eml} via l'interface web (client) et clique sur "analyser".
2. le client envoie une requête \texttt{http post} contenant le fichier multipart à la vue correspondante (\texttt{upload\_email()}) dans django.
3. la vue intercepte le fichier sous forme de flux mémoire et le transmet au module utilitaire (\texttt{parse\_eml\_in\_memory()}).
4. le parseur déstructure le mime, extrait le texte, le nettoie et isole les url, puis renvoie un dictionnaire python propre à la vue.
5. la vue transmet les données textuelles au module de machine learning encapsulé (\texttt{\_predict()}).
6. le modèle applique sa transformation tf-idf, exécute l'algorithme random forest, et retourne un score de probabilité.
7. la vue instancie un nouvel objet du modèle (\texttt{scanreport}) contenant le score et les métadonnées pour sauvegarde en base de données.
8. la vue déclenche une redirection vers le gabarit de résultat (\texttt{report\_detail.html}), qui affiche les données à l'utilisateur.
9. *post-traitement :* le fichier brut en mémoire est définitivement détruit, assurant la conformité au principe de "zero retention".

\section{modélisation de la base de données}

\subsection{choix du système de gestion de base de données (sgbd)}
bien que django soit configuré par défaut pour utiliser sqlite (idéal pour le prototypage rapide), le choix architectural final pour fakefinder s'est porté sur **postgresql**.
postgresql est un système de gestion de base de données relationnelle objet (sgbdro) open-source de niveau entreprise. son choix se justifie par plusieurs atouts critiques pour notre projet :
\begin{itemize}
    \item **support natif et avancé du json :** notre système doit stocker des listes de taille variable (comme les url extraites d'un e-mail) sans avoir à créer de tables relationnelles complexes et lourdes. postgresql gère les champs \texttt{jsonfield} de manière native, permettant d'indexer et de requêter directement à l'intérieur des structures json.
    \item **concurrence et performance :** contrairement à sqlite qui verrouille l'ensemble de la base lors d'une écriture, postgresql gère la concurrence multiversion (mvcc), permettant à de multiples utilisateurs d'enregistrer leurs rapports d'analyse simultanément sans dégradation des performances.
\end{itemize}

\subsection{diagramme de classes et structure des entités}
l'architecture de la persistance repose sur deux entités majeures : \texttt{scanreport} et \texttt{signature}, liées au modèle d'authentification utilisateur (\texttt{user}) fourni par défaut par django.

\begin{figure}[h]
    \centering
    \includegraphics[width=0.85\textwidth]{classes_diagram.png}
    \caption{diagramme de classes détaillant les entités du domaine, incluant le nouveau système de signatures}
\end{figure}

les attributs de la classe \texttt{scanreport} traduisent les exigences d'information :
\begin{itemize}
    \item \texttt{risk\_score} : un attribut de type chaîne de caractères restreint par des choix prédéfinis (\texttt{low}, \texttt{medium}, \texttt{high}), servant de base rapide pour le rendu visuel (couleurs) de l'interface.
    \item \texttt{score} : un entier (\texttt{integerfield}) compris entre 0 et 100. il représente la confiance granulaire du modèle d'intelligence artificielle dans sa détection de fraude.
    \item \texttt{suspicious\_urls} : un attribut de type \texttt{jsonfield} contenant un tableau des chaînes de caractères correspondant aux url isolées dans le message.
    \item \texttt{header\_anomalies} : un second \texttt{jsonfield} permettant de stocker sous forme de dictionnaire clé-valeur les métadonnées intéressantes (domaine de l'expéditeur, objet de l'e-mail).
\end{itemize}

\textbf{l'entité signature et domaines de confiance (nouveauté implementation) :}
pour optimiser les performances et permettre une détection immédiate des menaces connues (blacklisting dynamique) tout en réduisant les faux positifs, deux nouvelles entités ont été implémentées :
\begin{itemize}
    \item \texttt{signature} : stocke les empreintes cryptographiques (hashs) des indicateurs de compromission (urls, corps d'e-mail).
    \item \texttt{trusteddomain} : gère dynamiquement les domaines légitimes (ex: codeberg.org, github.com) pour appliquer des corrections de score intelligentes via l'interface d'administration.
\end{itemize}

\begin{itemize}
    \item \texttt{indicator\_hash} : stocke le hash sha-256 (pour les url) ou md5 (pour les corps d'e-mails).
    \item \texttt{indicator\_type} : définit la nature de la signature (\texttt{url}, \texttt{md5}, \texttt{sha256}).
    \item \texttt{description} : permet de documenter l'origine de la menace détectée.
\end{itemize}
lors d'un scan, le système interroge d'abord cette table (chemin court). en cas de correspondance, le verdict est instantanément classé comme "high risk" (score 100), évitant ainsi le coût de calcul de l'inférence ia pour les menaces déjà identifiées.
la cardinalité entre l'entité \texttt{user} et \texttt{scanreport} est de type "un-à-plusieurs" (1..n). cependant, l'utilisation de \texttt{on\_delete=models.set\_null} permet au système d'accepter des rapports d'analyse de manière anonyme (où la clé étrangère vers l'utilisateur reste nulle).


% ============================================================
%  chapitre 4 : réalisation et développement
% ============================================================
\chapter{réalisation technique et modélisation ia}

le chapitre de réalisation détaille l'implémentation concrète des concepts d'ingénierie logicielle et de data science évoqués lors de la conception. il s'agit du cœur technique du projet fakefinder.

\section{le moteur d'extraction et de traitement (parsing)}
avant même de faire intervenir l'intelligence artificielle, le système doit être capable de "lire" l'e-mail soumis de manière sécurisée. la tâche est complexe car le format \texttt{.eml} (norme mime) peut encapsuler des données textuelles, des images, des pièces jointes, et de multiples encodages de caractères au sein d'un même fichier.

le module utilitaire \texttt{utils.py} a été développé avec une logique de parsing défensive :
\begin{itemize}
    \item **validation préliminaire :** le script vérifie la taille du flux en mémoire. si celle-ci excède 10 mo, la fonction lève immédiatement une exception, protégeant le serveur contre l'épuisement de la ram.
    \item **navigation dans l'arborescence mime :** en utilisant la bibliothèque standard python \texttt{email}, le script parcourt itérativement (\texttt{msg.walk()}) toutes les sous-parties de l'e-mail. il ignore systématiquement les pièces jointes (ex: fichiers pdf ou images encapsulées en base64) pour se concentrer exclusivement sur les parties de type \texttt{text/plain} et \texttt{text/html}.
    \item **gestion de l'obfuscation d'encodage :** les attaquants manipulent souvent les déclarations de charset (jeux de caractères) pour perturber l'analyse. le script tente de décoder le bloc de texte avec le charset déclaré. si cela échoue (erreur \texttt{lookuperror}), il bascule automatiquement sur un forçage en \texttt{utf-8} avec l'option \texttt{errors="ignore"}, garantissant la récupération du texte, même si certains caractères spéciaux sont perdus.
    \item **expressions régulières (regex) :** une fois le texte nettoyé des espaces excédentaires, le script exécute une regex robuste compilée (\texttt{r'(https?://[\^\textbackslash s>"'\textbackslash ']+)'}) pour isoler méthodiquement tous les liens hyperboliques contenus dans le message.
\end{itemize}

\section{le modèle de machine learning : traitement et entraînement}
le succès d'un système de détection de phishing repose sur la qualité de son modèle prédictif. nous avons opté pour une approche d'apprentissage supervisé, ce qui implique que le modèle "apprend" à différencier le bien du mal à partir d'exemples d'e-mails préalablement labellisés.

\subsection{acquisition des données et limitations}
le modèle a été entraîné sur le dataset public \textit{"phishing email detection"} publié par subhalaxmi rout sur la plateforme kaggle. ce jeu de données rassemble environ 18 650 e-mails, offrant une répartition relativement équilibrée : $\sim$56\% d'e-mails légitimes (safe) et $\sim$44\% d'e-mails frauduleux (phishing).

le processus de préparation des données a inclus :
\begin{itemize}
    \item \textbf{nettoyage :} suppression des entrées vides, des doublons exacts et normalisation des labels (conversion des chaînes "phishing email" et "safe email" en valeurs binaires 1 et 0).
    \item \textbf{analyse des classes :} avec 10 320 e-mails "safe" et 8 330 e-mails "phishing", le dataset est suffisamment équilibré pour éviter un biais majeur sans nécessiter de techniques de suréchantillonnage (smote).
    \item \textbf{limitations :} la quasi-totalité des e-mails y figure en langue anglaise. par conséquent, bien que l'analyse structurelle reste valide, la détection sémantique sur des campagnes de phishing ciblant spécifiquement d'autres langues sera statistiquement moins efficace. de plus, le dataset datant de quelques années, il ne contient pas nécessairement les schémas d'attaques générées par les ia les plus récentes.
\end{itemize}

\subsection{étude comparative et choix du modèle}
avant de fixer l'architecture prédictive sur les forêts aléatoires, une phase d'expérimentation a permis de comparer plusieurs algorithmes de classification :

\begin{table}[h]
\centering
\caption{comparaison des performances des algorithmes (sur le jeu de test initial)}
\begin{tabular}{lcccc}
\toprule
\textbf{algorithme} & \textbf{exactitude} & \textbf{précision} & \textbf{rappel} & \textbf{temps inférence} \\
\midrule
régression logistique & 91,2\% & 89,5\% & 88,2\% & < 10ms \\
naive bayes (multinomial) & 88,4\% & 85,1\% & 92,4\% & < 5ms \\
svm (noyau linéaire) & 95,1\% & 93,2\% & 94,8\% & ~150ms \\
\textbf{random forest} & \textbf{97,0\%} & \textbf{94,0\%} & \textbf{97,0\%} & \textbf{~40ms} \\
\bottomrule
\end{tabular}
\end{table}

\textbf{justification du choix :}
\begin{itemize}
    \item \textit{régression logistique :} trop simpliste pour capter les interactions complexes entre les termes.
    \item \textit{naive bayes :} performant en rappel mais génère trop de faux positifs à cause de l'hypothèse d'indépendance des mots.
    \item \textit{svm :} très précis, mais le coût computationnel et la difficulté d'obtenir des probabilités calibrées (nécessitant la méthode de platt) le rendent moins pratique que random forest.
    \item \textbf{random forest :} offre le meilleur compromis entre robustesse, précision et rapidité. il gère naturellement les données de haute dimension issues du tf-idf et permet une interprétabilité via l'importance des variables.
\end{itemize}

\subsection{l'ingénierie des caractéristiques (feature engineering)}
les algorithmes d'apprentissage automatique ne comprennent que des matrices numériques. il est donc indispensable de convertir le texte brut des e-mails en représentations mathématiques. fakefinder combine deux approches complémentaires au sein d'un \texttt{combinedfeaturetransformer}.

\textbf{1. les caractéristiques structurelles (heuristique) :}
un module python personnalisé (\texttt{structuralfeatureextractor}) parcourt le texte pour en tirer des statistiques d'ordre dimensionnel et émotionnel. dix caractéristiques numériques (features) sont calculées, parmi lesquelles :
\begin{itemize}
    \item \texttt{url\_count} et \texttt{html\_link\_count} : le nombre total de liens détectés et de balises \texttt{href} (les e-mails de phishing contiennent souvent une densité anormale de liens pour maximiser les chances de clic).
    \item \texttt{has_short_url} : un booléen (0 ou 1) indiquant la présence de domaines connus pour raccourcir les url (\texttt{bit.ly}, \texttt{tinyurl.com}, \texttt{t.co}, etc.).
    \item \texttt{phishing_keyword_score} : un compteur mesurant la fréquence d'apparition d'un vocabulaire hautement suspect (plus de 40 mots-clés tels que "verify", "suspended", "password", "bank", "alert").
    \item \texttt{dollar_count} et \texttt{exclamation_count} : le nombre de signes monétaires et de points d'exclamation, souvent corrélés aux fausses promesses de gains ou aux alertes graves.
    \item \texttt{has_html_tags} et \texttt{has_email_in_body} : détection de la structure html et de la présence d'adresses e-mails dans le corps, indicateurs de messages formatés ou de sollicitations directes.
    \item \texttt{normalized_length} et \texttt{caps_word_count} : mesure de la longueur du message et de l'utilisation excessive de majuscules (cri numérique).
\end{itemize}

\textbf{2. la vectorisation sémantique tf-idf :}
pour capturer le sens global du message au-delà des mots-clés prédéfinis, nous avons implémenté l'algorithme \textbf{tf-idf} (term frequency - inverse document frequency). le tf-idf évalue l'importance d'un mot (ou d'une combinaison de deux mots consécutifs, les bigrammes) dans un document donné, par rapport à l'ensemble du corpus d'entraînement.

la pondération tf-idf d'un terme $i$ dans un document $j$ est définie par l'équation suivante :
\begin{equation}
    w_{i,j} = tf_{i,j} \times \log\left(\frac{n}{df_i}\right)
\end{equation}
où :
\begin{itemize}
    \item $tf_{i,j}$ (term frequency) représente la fréquence d'apparition brute du terme $i$ dans l'e-mail $j$.
    \item $df_i$ (document frequency) représente le nombre total d'e-mails dans le corpus contenant le terme $i$.
    \item $n$ est le nombre total de documents dans le dataset d'entraînement.
\end{itemize}
la force du tf-idf, imposée par l'utilisation du logarithme inverse, est qu'il pénalise fortement les mots très fréquents dans la langue (les mots vides ou "stop words" tels que "le", "de", "and") tout en accordant un poids mathématique prépondérant aux mots plus rares et sémantiquement porteurs de sens (ex: "invoice", "validate", "account"). le \texttt{tfidfvectorizer} configuré génère ainsi une matrice creuse monumentale (sparse matrix) comportant jusqu'à 20 000 dimensions (colonnes) par e-mail.

\subsection{justification et implémentation de l'algorithme : random forest}
avant de fixer l'architecture prédictive, une étude comparative des modèles de classification classiques a été réalisée. 
la \textit{régression logistique} s'est avérée trop simpliste pour capter les interactions non-linéaires complexes entre les termes du vocabulaire. le \textit{naive bayes}, souvent utilisé en filtrage de spam pour sa rapidité, part du principe que toutes les caractéristiques (mots) sont indépendantes les unes des autres, ce qui est faux dans le contexte d'une phrase manipulatrice. les \textit{machines à vecteurs de support (svm)}, bien que très performantes en haute dimension, nécessitaient des temps de calcul et d'optimisation prohibitifs.

notre choix final s'est porté sur l'algorithme des \textbf{forêts aléatoires (random forest)}. introduit par leo breiman en 2001, il repose sur le principe de l'apprentissage d'ensemble (ensemble learning), et plus particulièrement sur la technique du \textit{bagging} (bootstrap aggregating).

**principe de fonctionnement :**
le random forest ne repose pas sur un seul arbre de décision complexe et profond, qui aurait une fâcheuse tendance à apprendre le jeu de données par cœur (surapprentissage ou overfitting). au contraire, lors de la phase d'entraînement, l'algorithme génère une multitude d'arbres de décision indépendants (dans notre cas, 200 arbres). 
chaque arbre est entraîné sur un sous-échantillon aléatoire (tirage avec remise) des données d'entraînement, et n'a accès qu'à un sous-ensemble aléatoire des caractéristiques (dimensions tf-idf) à chaque nœud pour décider de sa coupure, généralement basée sur la minimisation de l'\textbf{impureté de gini} :
\begin{equation}
    gini(p) = 1 - \sum_{i=1}^{c} p_i^2
\end{equation}
où $p_i$ est la probabilité qu'un élément appartienne à la classe $i$ (safe ou phishing).

lors de la phase de prédiction d'un nouvel e-mail, chaque arbre de la forêt "vote" indépendamment pour une classe. le résultat final rendu par le modèle est la moyenne ou le mode de tous ces votes, garantissant une prédiction robuste, stable et beaucoup moins sensible aux variations ou aux bruits dans le texte de l'e-mail.

**calibration des probabilités :**
les random forests ont tendance à générer des probabilités biaisées (les scores se regroupent souvent autour de la médiane de 50\%). pour pallier cela et offrir un score de 0 à 100 fiable à l'utilisateur final, nous avons encapsulé le modèle dans un \texttt{calibratedclassifiercv} utilisant la régression isotonique. ce processus de calibration permet de transformer les votes de la forêt en probabilités réelles, alignées sur la fréquence d'apparition des classes dans le jeu de test. le pipeline complet a ensuite été exporté sur disque sous la forme d'un fichier binaire \texttt{model.joblib} pesant plusieurs dizaines de mégaoctets.

\section{développement de l'interface utilisateur (frontend)}
le défi du développement frontend consistait à masquer la lourdeur et la complexité des opérations mathématiques sous-jacentes derrière une interface délibérément minimaliste, propre et intuitive. l'application doit rassurer l'utilisateur par son design moderne.

\begin{figure}[h]
    \centering
    \includegraphics[width=0.8\textwidth]{frontend_0.jpeg}
    \caption{interface de la page d'accueil : zone d'upload de type "drag and drop"}
\end{figure}

l'interaction principale s'effectue via une zone de dépôt centralisée (drag and drop) gérée par un script javascript. ce dernier assure un premier niveau de validation côté client en vérifiant l'extension du fichier (\texttt{.eml}) avant même de solliciter le serveur, économisant ainsi de la bande passante et améliorant la réactivité perçue.

\begin{figure}[h]
    \centering
    \includegraphics[width=0.8\textwidth]{frontend_1.jpeg}
    \caption{tableau de bord de restitution : présentation des résultats d'analyse et du verdict ia}
\end{figure}

une fois l'analyse terminée (processus synchrone prenant moins d'une seconde en production), l'utilisateur est redirigé vers le tableau de bord de résultats. l'ergonomie repose sur un retour visuel instantané : une jauge et un code couleur sémantique (vert = sûr, orange = suspect, rouge = dangereux) encadrent le score de confiance de l'algorithme.
dans un souci de transparence et de pédagogie, le rapport détaille également les éléments techniques factuels ayant conduit à ce score : le domaine réel d'expédition (souvent très différent du nom d'affichage) et la liste exhaustive des liens extraits, permettant à l'utilisateur d'exercer son propre esprit critique.


% ============================================================
%  chapitre 5 : perspectives technologiques
% ============================================================
\chapter{perspectives technologiques et systèmes hybrides}

\textit{note sur l'évolution du système : alors que la version initiale de fakefinder reposait exclusivement sur l'apprentissage automatique, l'implémentation a évolué vers une architecture hybride. le système intègre désormais une base de signatures locales (v1.1) comme socle de détection, tout en conservant l'api virustotal comme perspective majeure d'évolution logicielle.}

\section{l'approche hybride : signatures locales et ia}
le modèle prédictif développé (random forest + tf-idf) a démontré d'excellentes capacités pour l'analyse sémantique. néanmoins, l'ia est coûteuse en ressources de calcul et peut parfois être sujette à l'incertitude. pour optimiser le pipeline, nous avons implémenté un système de \textbf{signatures locales}.

cette couche de détection agit comme un "pare-feu applicatif" de premier niveau :
\begin{itemize}
    \item \textbf{chemin court (fast-path) :} avant d'invoquer le modèle machine learning, le système calcule le hash (sha-256) des url et le hash (md5) du corps du message.
    \item \textbf{lookup immédiat :} ces empreintes sont comparées en temps réel avec une table de \texttt{signatures} en base de données.
    \item \textbf{verdict instantané :} si une correspondance est trouvée (menace déjà connue et répertoriée par l'administrateur), le scan s'arrête et renvoie un score de 100\% (phishing). cette approche garantit une performance optimale et une fiabilité absolue sur les menaces récurrentes.
\end{itemize}

\section{la perspective virustotal (threat intelligence)}
malgré l'ajout des signatures locales, la détection reste limitée aux données connues par notre système propre. pour pallier cette limite, la stratégie d'évolution prévoit l'intégration d'un composant de \textbf{cyber threat intelligence (cti)} externe via l'api virustotal.

\section{architecture de l'intégration virustotal prévue}
l'intégration de l'api rest de virustotal (filiale de google) permettra d'agréger les résultats d'analyse de plus de 70 moteurs de sécurité mondiaux. le workflow technique futur s'articulera selon les étapes suivantes :
\begin{enumerate}
    \item \textbf{extraction syntaxique :} les url sont extraites via les expressions régulières lors du parsing.
    \item \textbf{requêtage asynchrone :} le backend interrogera l'api virustotal de manière non-bloquante.
    \item \textbf{pondération hybride :} le score final sera une fusion entre l'ia locale, les signatures internes et la réputation mondiale virustotal.
\end{enumerate}

cette évolution transformera fakefinder en un véritable agrégateur hybride d'investigation de sécurité, comparable aux outils de pointe des centres opérationnels de sécurité (soc).

cette évolution majeure, prévue pour la prochaine itération, transformera fakefinder. de simple filtre de classification local, il deviendra un véritable agrégateur hybride d'investigation de sécurité, comparable aux outils déployés dans les centres opérationnels de sécurité (soc).


% ============================================================
%  chapitre 6 : tests, validation et conclusion
% ============================================================
\chapter{évaluation, tests et conclusion générale}

la phase de test et de validation est cruciale en génie logiciel, et d'autant plus dans le domaine de la cybersécurité où les erreurs de jugement (faux négatifs) peuvent avoir des conséquences désastreuses. l'évaluation de fakefinder a été menée sur deux axes distincts : la validation statistique de son "cerveau" (le modèle machine learning) et les tests fonctionnels de son enveloppe (l'application web django).

\section{validation du modèle de machine learning}
la méthodologie d'évaluation d'un modèle prédictif impose de ne jamais le tester sur les données qu'il a utilisées pour s'entraîner. sur le dataset total de $\sim$18 650 e-mails, 20\% des données (soit environ 3 730 e-mails) ont été formellement isolées pour constituer le "jeu de test" (test set).

\subsection{métriques d'évaluation théoriques}
la performance a été mesurée à l'aide des métriques standards de classification, dérivées des concepts de vrais positifs (vp), faux positifs (fp), vrais négatifs (vn) et faux négatifs (fn) :
\begin{itemize}
    \item \textbf{précision (precision) :} mesure la proportion d'identifications correctes parmi tous les e-mails classés comme phishing par le modèle. formule : $\frac{vp}{vp + fp}$.
    \item \textbf{rappel (recall / sensibilité) :} mesure la proportion d'e-mails de phishing réels qui ont été correctement détectés par le modèle. formule : $\frac{vp}{vp + fn}$.
    \item \textbf{exactitude globale (accuracy) :} le pourcentage total de prédictions correctes sur l'ensemble du jeu de test. formule : $\frac{vp + vn}{vp + vn + fp + fn}$.
\end{itemize}

\subsection{résultats statistiques obtenus}
le rapport de classification généré par `scikit-learn` sur le jeu de test a révélé des performances dépassant nos attentes initiales pour une application en environnement de production :

\begin{table}[h]
\centering
\caption{rapport de classification détaillé du modèle random forest sur le jeu de test}
\begin{tabular}{lrrrr}
\toprule
\textbf{classe prédite} & \textbf{précision} & \textbf{rappel} & \textbf{score f1} & \textbf{volume (support)} \\
\midrule
emails légitimes (safe)  & 0,98 & 0,96 & 0,97 & 2 265 \\
emails frauduleux (phish)& 0,94 & 0,97 & 0,96 & 1 462 \\
\midrule
\textbf{exactitude globale (accuracy)} & \multicolumn{4}{r}{\textbf{97\%}} \\
\textbf{score roc-auc} & \multicolumn{4}{r}{\textbf{0,9928}} \\
\bottomrule
\end{tabular}
\end{table}

\textbf{analyse et interprétation :}
avec une **précision de 94\%** sur la classe "phishing", cela signifie que lorsque fakefinder alerte l'utilisateur en déclarant un e-mail malveillant, il a mathématiquement raison dans 94\% des cas. le taux de faux positifs est donc restreint, ce qui évite la "fatigue d'alerte" chez l'utilisateur (le fait de finir par ignorer un outil de sécurité qui sonne trop souvent pour des e-mails normaux).
plus important encore en cybersécurité, le **rappel (recall) s'élève à 97\%**. le modèle est donc extrêmement sensible aux menaces : il parvient à débusquer et intercepter 97\% des e-mails frauduleux présents dans l'échantillon, minimisant ainsi le risque critique qu'une attaque ne passe à travers les mailles du filet (faux négatifs).
le score **roc-auc de 0,9928** (très proche de la perfection qui est 1.0) confirme que le modèle possède une capacité de discrimination exceptionnelle entre les deux classes.

\begin{figure}[h]
    \centering
    \includegraphics[width=0.6\textwidth]{image_matrice_confusion.png}
    \caption{représentation visuelle de la matrice de confusion (confusion matrix)}
\end{figure}

\section{tests fonctionnels et de robustesse de l'application web}
au-delà de la data science, le système logiciel applicatif a été éprouvé dans des conditions d'utilisation imitant un environnement réel de production, afin de valider la conformité aux exigences fonctionnelles (fr) et non fonctionnelles (nfr).

\begin{table}[h]
\centering
\caption{matrice d'exécution des scénarios de tests fonctionnels web}
\resizebox{\textwidth}{!}{%
\begin{tabular}{|l|p{6cm}|p{6cm}|c|}
\hline
\textbf{id test} & \textbf{description du scénario (input)} & \textbf{comportement système attendu} & \textbf{statut final} \\ \hline
tc\_web\_01 & soumission via le formulaire d'un fichier \texttt{.eml} valide, bien formaté et de taille standard. & succès de l'upload, traitement synchrone, redirection vers la page du rapport généré affichant le score calculé. & \checkmark succès \\ \hline
tc\_web\_02 & tentative d'upload d'un fichier invalide ou potentiellement dangereux (ex: un script \texttt{.py} ou un \texttt{.exe}). & rejet natif par le validateur de formulaire django avant tout traitement. message d'erreur clair affiché à l'utilisateur. & \checkmark succès \\ \hline
tc\_web\_03 & tentative de soumission d'un fichier extraordinairement volumineux (fichier > 10 mo) simulant une attaque dos. & rejet automatique par les sécurités du middleware et de la vue. renvoi d'un code http 400 (bad request). & \checkmark succès \\ \hline
tc\_web\_04 & soumission d'un fichier \texttt{.eml} volontairement corrompu (sans en-têtes ou corps de texte vide). & le parseur lève et gère l'exception sans faire crasher le serveur (pas de page blanche http 500). un score neutre par défaut est attribué. & \checkmark succès \\ \hline
tc\_web\_05 & soumission d'un e-mail ne contenant aucun lien (url) dans son corps. & le système doit traiter l'e-mail normalement, le score se basant alors uniquement sur l'analyse sémantique et les autres caractéristiques structurelles. & \checkmark succès \\ \hline
tc\_web\_06 & utilisateur visiteur (non authentifié) effectuant une analyse d'e-mail. & le processus d'analyse réussit. le rapport est affiché mais n'est rattaché à aucun profil dans la base de données (historique vierge). & \checkmark succès \\ \hline
tc\_web\_07 & utilisateur inscrit et authentifié effectuant une analyse d'e-mail. & l'analyse réussit. les métadonnées du rapport sont sauvegardées en bdd et attachées à l'id de l'utilisateur. il est visible dans le tableau de bord. & \checkmark succès \\ \hline
tc\_web\_08 & test d'isolation des données : utilisateur "a" tente de lire l'url du rapport d'analyse de l'utilisateur "b". & le système identifie que le rapport appartient à un tiers. la vue bloque l'accès et renvoie un code http 403 (accès interdit / forbidden). & \checkmark succès \\ \hline
tc\_web\_09 & tentative de suppression d'un rapport par un utilisateur authentifié qui n'en est pas le propriétaire. & la requête de suppression est rejetée avec un code http 403. & \checkmark succès \\ \hline
tc\_web\_10 & gestion des erreurs lors d'une tentative de connexion avec des identifiants erronés. & message d'erreur clair affiché ("identifiant ou mot de passe incorrect") et incrémentation du compteur de tentatives pour protection anti-bruteforce. & \checkmark succès \\ \hline
\end{tabular}%
}
\end{table}

\subsection{évaluation des performances temporelles}
l'un des défis de l'architecture consistait à garantir une latence minimale. le traitement complet de bout-en-bout (upload http du fichier, parsing en mémoire vive de l'arborescence mime complexe, extraction et tokenization du texte, inférence de l'algorithme random forest, sauvegarde postgresql des métadonnées, et rendu de la page html) a été chronométré.
en environnement de développement local standard, cette boucle complète s'exécute en **450 millisecondes en moyenne**. cette vélocité, rendue possible par le "lazy loading" du modèle ml et le traitement in-memory (évitant les goulots d'étranglement des lectures/écritures sur disque), respecte très largement l'exigence initiale fixée à moins de 5 secondes.

\section{conclusion générale et bilan du projet}
l'accomplissement du projet de fin d'année \textbf{fakefinder} s'est soldé par la création d'une solution logicielle robuste, moderne et éthique, capable de répondre efficacement à l'une des préoccupations majeures de la cybersécurité contemporaine : la recrudescence endémique des campagnes de hameçonnage par courrier électronique.

l'objectif principal était de fournir un outil d'investigation automatisé accessible, sans obliger l'utilisateur à sacrifier la confidentialité de ses correspondances. cet objectif a été pleinement atteint. en associant la structure éprouvée, sécurisée et rapide du framework web python django, à la puissance de classification et d'analyse sémantique offerte par les algorithmes ensemblistes (random forest) de \texttt{scikit-learn}, nous avons démontré la viabilité technique d'un moteur d'analyse délégué au serveur. l'architecture respecte scrupuleusement le principe fondateur du "zero retention" : aucune ligne de l'e-mail brut n'est archivée par le système.

sur le plan purement pédagogique, la réalisation de ce projet sur une durée condensée de 5 semaines a représenté un défi particulièrement stimulant. il nous a offert une opportunité précieuse d'explorer, de maîtriser et de consolider des concepts transversaux hautement valorisés dans l'ingénierie informatique moderne :
\begin{itemize}
    \item **le génie logiciel avancé :** implémentation stricte du paradigme d'architecture mvt, maîtrise d'un orm relationnel, respect du cycle de vie du développement logiciel (sdlc) via les méthodes agiles, et sécurisation applicative défensive contre les failles web courantes.
    \item **la data science et l'intelligence artificielle :** expérimentation complète du pipeline de machine learning (acquisition de données, analyse exploratoire, ingénierie des caractéristiques textuelles via tf-idf, entraînement, sérialisation de modèles, et évaluation scientifique par métriques statistiques avancées).
    \item **la cybersécurité :** compréhension de l'ingénierie sociale, des techniques d'obfuscation des attaquants et de l'architecture des protocoles de messagerie.
\end{itemize}

les résultats statistiques (précision de 94\% et rappel de 97\%) attestent de la fiabilité immédiate de la solution sur les corpus anglophones, confirmant que l'analyse algorithmique des motifs langagiers et structurels supplante largement l'approche réactive des simples listes noires.

enfin, la clarté de la base architecturale (code source modulaire) ouvre sereinement la voie aux évolutions et améliorations futures. comme détaillé dans nos perspectives technologiques, l'intégration à venir d'interfaces de programmation (api) de \textit{cyber threat intelligence}, telles que virustotal, viendra parachever ce système. cette synergie entre une analyse heuristique sémantique locale (ia) et une vérification factuelle et globale de la réputation des liens (cloud) transformera définitivement fakefinder en un véritable dispositif hybride de diagnostic de menaces, prêt à affronter les défis complexes de la cybersécurité de demain.

% ============================================================
%  bibliographie (norme apa stricte - exigence emsi)
% ============================================================
\chapter*{bibliographie}
\addcontentsline{toc}{chapter}{bibliographie}

\begin{itemize}[leftmargin=0.5cm, itemindent=-0.5cm, label={}]
  
  \item abu-nimeh, s., nappa, d., wang, x., \& nair, s. (2007). a comparison of machine learning techniques for phishing detection. \textit{proceedings of the anti-phishing working groups 2nd annual ecrime researchers summit}, 60-69.
  
  \item basnet, r., sung, a. h., \& liu, q. (2008). rule-based phishing attack detection. \textit{international conference on security and management}, 22-28.
  
  \item breiman, l. (2001). random forests. \textit{machine learning}, 45(1), 5-32.
  
  \item django software foundation. (2024). \textit{django documentation - version 6.0}. repéré à https://docs.djangoproject.com/
  
  \item hannousse, a., \& yahiouche, s. (2021). towards benchmark datasets for machine learning based website phishing detection: an experimental study. \textit{engineering applications of artificial intelligence}, 104, 104347.
  
  \item mohammad, r. m., thabtah, f., \& mccluskey, l. (2014). predicting phishing websites based on self-structuring neural network. \textit{neural computing and applications}, 25(2), 443-458.
  
  \item pedregosa, f., varoquaux, g., gramfort, a., michel, v., thirion, b., grisel, o., ... \& duchesnay, e. (2011). scikit-learn: machine learning in python. \textit{journal of machine learning research}, 12, 2825-2830.
  
  \item rout, s. (2022). \textit{phishing email detection dataset}. kaggle. repéré à https://www.kaggle.com/datasets/subhajournal/phishingemails
  
  \item sahingoz, o. k., batur, e., bulut, f., \& taskin, a. (2019). machine learning based phishing detection from urls. \textit{expert systems with applications}, 117, 345-357.
  
  \item verizon. (2024). \textit{data breach investigations report (dbir) 2024}. repéré à https://www.verizon.com/business/resources/reports/dbir/
  
\end{itemize}

% ============================================================
%  annexes
% ============================================================
\appendix
\chapter{listings complets du code source}

\section{Module d'Extraction des Caractéristiques (ml/features.py)}
\begin{lstlisting}[language=Python]
"""
Feature extraction for phishing email detection.
Combines TF-IDF on full text with 10 structural features.
"""
import re
import numpy as np
import scipy.sparse as sp
from sklearn.base import BaseEstimator, TransformerMixin
from sklearn.feature_extraction.text import TfidfVectorizer

PHISHING_KEYWORDS = [
    "verify", "account", "update", "confirm", "click here", "urgent",
    "immediately", "suspended", "limited", "expires", "login", "password",
    "bank", "credit", "security", "alert", "prize", "validate",
]
SHORT_URL_PATTERN = re.compile(r"bit\.ly|tinyurl\.com|goo\.gl|t\.co|rb\.gy", re.IGNORECASE)
URL_PATTERN = re.compile(r"https?://[^\s<>\"']+", re.IGNORECASE)

class StructuralFeatureExtractor(BaseEstimator, TransformerMixin):
    def fit(self, X, y=None): return self
    def transform(self, X):
        rows = []
        for text in X:
            t = text.lower()
            words = text.split()
            urls = URL_PATTERN.findall(t)
            html_links = len(re.findall(r"href=", t))
            rows.append([
                len(urls), html_links,
                int(bool(SHORT_URL_PATTERN.search(t))),
                min(text.count("!"), 20), min(text.count("$"), 10),
                sum(1 for kw in PHISHING_KEYWORDS if kw in t),
                int("<html" in t or "<div" in t),
                int(bool(re.search(r"[\w.+-]+@[\w-]+\.[a-z]{2,}", t))),
                min(len(words), 3000) / 3000,
                sum(1 for w in words[:60] if len(w) > 2 and w.isupper())
            ])
        return sp.csr_matrix(np.array(rows, dtype=float))

class CombinedFeatureTransformer:
    def __init__(self, max_features=20000):
        self.tfidf = TfidfVectorizer(max_features=max_features, ngram_range=(1, 2))
        self.structural = StructuralFeatureExtractor()
    def fit(self, X, y=None):
        self.tfidf.fit(X)
        return self
    def transform(self, X):
        return sp.hstack([self.tfidf.transform(X), self.structural.transform(X)])

def prepare_email_text(parsed):
    parts = [parsed.get("subject", ""), parsed.get("body_text", ""), 
             " ".join(parsed.get("urls", []))]
    return " ".join(p for p in parts if p)
\end{lstlisting}

\newpage
\section{Script d'Entraînement de l'Intelligence Artificielle (ml/train.py)}
\begin{lstlisting}[language=Python]
"""
Entrainement du modele Random Forest avec calibration.
Dataset: Phishing_Email.csv (Kaggle).
"""
import joblib
import pandas as pd
from sklearn.ensemble import RandomForestClassifier
from sklearn.calibration import CalibratedClassifierCV
from sklearn.model_selection import train_test_split
from ml.features import CombinedFeatureTransformer

def train(csv_path, output_path):
    df = pd.read_csv(csv_path).dropna()
    X = df["Email Text"].values
    y = df["Email Type"].apply(lambda x: 1 if "Phish" in str(x) else 0).values

    X_train, X_test, y_train, y_test = train_test_split(X, y, test_size=0.2, stratify=y)
    
    transformer = CombinedFeatureTransformer()
    X_train_feat = transformer.fit_transform(X_train)
    
    # RandomForest + Isotonic Calibration pour des probabilites fiables
    base_clf = RandomForestClassifier(n_estimators=200, class_weight="balanced")
    clf = CalibratedClassifierCV(base_clf, method="isotonic", cv=5)
    clf.fit(X_train_feat, y_train)

    joblib.dump({"transformer": transformer, "classifier": clf}, output_path)
\end{lstlisting}

\newpage
\section{Moteur de Parsing d'E-mails en Mémoire Vive (scanner/utils.py)}
\begin{lstlisting}[language=Python]
"""
Moteur Zero Retention pour la manipulation des fichiers .eml.
Extraction des metadonnees et du texte sans persistance disque.
"""
import email
import re

URL_PATTERN = re.compile(r'(https?://[^\s>"\']+)')

def parse_eml_in_memory(file_obj):
    file_obj.seek(0)
    msg = email.message_from_bytes(file_obj.read())
    body_text = ""
    
    if msg.is_multipart():
        for part in msg.walk():
            if part.get_content_type() in ["text/plain", "text/html"]:
                payload = part.get_payload(decode=True)
                if payload: body_text += payload.decode(errors="ignore") + " "
    else:
        payload = msg.get_payload(decode=True)
        if payload: body_text = payload.decode(errors="ignore")
            
    body_text = re.sub(r"\s+", " ", body_text).strip()
    urls = list(set(URL_PATTERN.findall(body_text)))
    
    sender = str(msg.get("From", ""))
    email_match = re.search(r"[\w.+-]+@[\w-]+\.[a-z]{2,}", sender)
    sender_email = email_match.group(0) if email_match else ""
    
    return {
        "sender_email": sender_email,
        "sender_domain": sender_email.split("@")[-1] if "@" in sender_email else "",
        "subject": str(msg.get("Subject", "")),
        "body_text": body_text,
        "urls": urls
    }
\end{lstlisting}

\newpage
\section{Structure des Données et Signatures (scanner/models.py)}
\begin{lstlisting}[language=Python]
from django.db import models

class ScanReport(models.Model):
    risk_score = models.CharField(max_length=10, choices=[
        ("LOW", "Low (Green)"), ("MEDIUM", "Medium (Orange)"), ("HIGH", "High (Red)")
    ], default="LOW")
    score = models.IntegerField(default=0)
    suspicious_urls = models.JSONField(default=list, blank=True)
    header_anomalies = models.JSONField(default=list, blank=True)

class Signature(models.Model):
    """Stockage des hashs de phishing connus."""
    TYPE_CHOICES = [('URL', 'URL Hash'), ('MD5', 'File MD5')]
    indicator_hash = models.CharField(max_length=64, unique=True, db_index=True)
    indicator_type = models.CharField(max_length=10, choices=TYPE_CHOICES)
    description = models.CharField(max_length=255, blank=True)
\end{lstlisting}

\newpage
\section{Logique de Détection Hybride (scanner/views.py)}
\begin{lstlisting}[language=Python]
import hashlib
from .models import Signature

def _check_signatures(parsed_data):
    # Verification des URLs (SHA256)
    urls = parsed_data.get("urls", [])
    if urls:
        url_hashes = [hashlib.sha256(u.encode()).hexdigest() for u in urls]
        if Signature.objects.filter(indicator_type='URL', indicator_hash__in=url_hashes).exists():
            return "HIGH", 100

    # Verification du corps (MD5)
    body = (parsed_data.get("body_text", "")).encode()
    if body:
        md5_hash = hashlib.md5(body).hexdigest()
        if Signature.objects.filter(indicator_type='MD5', indicator_hash=md5_hash).exists():
            return "HIGH", 100
    return None

def _predict(parsed_data):
    # Chemin court : Signatures
    result = _check_signatures(parsed_data)
    if result: return result

    # Chemin long : Intelligence Artificielle
    # ... Inférence Random Forest ...
\end{lstlisting}

\end{document}
